% claim.tex -- tatami-mod8-defect priority claim note (SOP 08b claim channel).
% AI-generated 2026-09-27 by the AI4Math paper department (claim lane).
% lstlean.tex copied unchanged from harness/templates/paper/lstlean.tex.
% Machine-readable anchors: every theorem carries a % LEAN: line (task id +
% declaration name + grade); paper-lint.sh and the claim audit check against it.
% All Lean listings are byte-exact mechanical excerpts of the frozen final
% development tasks/20260926-beian-select/final/01-proof-complete.lean.
\documentclass[11pt]{article}

\usepackage[T1]{fontenc}
\usepackage[utf8]{inputenc}
\usepackage{amsmath,amssymb,amsthm}
\usepackage{hyperref}
% lstlean.tex — Lean style for the listings package.
% Lineage: Jeremy Avigad (2015), adapted from Assia Mahboubi's SSR style;
% inspired by Olivier Verdier's unixode.sty for the Unicode literate table.
% No CTAN package or upstream repo exists; papers carry this file in their
% source package (cap set arXiv:1907.01449, sphere eversion arXiv:2210.07746).
% AI4Math adaptation (AI-generated, 2026-09-26): sorry/admit/sorryAx elevated
% to a dedicated error tier, matching the pipeline's 0-sorry constitution.
\usepackage{listings}
\usepackage{xcolor}

\definecolor{keywordcolor}{rgb}{0.7, 0.1, 0.1}   % red keywords
\definecolor{tacticcolor}{rgb}{0.0, 0.1, 0.6}    % blue tactics
\definecolor{commentcolor}{rgb}{0.4, 0.4, 0.4}   % grey comments
\definecolor{symbolcolor}{rgb}{0.0, 0.1, 0.6}    % blue symbols
\definecolor{sortcolor}{rgb}{0.1, 0.5, 0.1}      % green sorts
\definecolor{errorcolor}{rgb}{0.6, 0, 0}         % dark red: sorry/admit
\definecolor{stringcolor}{rgb}{0.5, 0.1, 0.5}    % purple strings

\lstdefinelanguage{lean}{
  % Anything between $ becomes LaTeX math mode
  mathescape=false,
  % Comments may or not include Latex commands
  texcl=false,
  % keywords, list taken from lean-syntax.txt
  morekeywords=[1]{import, prelude, protected, private, noncomputable,
    definition, meta, renaming, hiding, parameter, parameters, begin,
    constant, constants, lemma, variable, variables, theory, print,
    theorem, example, open, section, namespace, universe, universes, end,
    modifier, modifiers, using, export, exposing, axiom, inductive,
    coinductive, with, structure, class, instance, attribute, local, where,
    by, have, show, let, in, if, then, else, match, do, for, fun, assume,
    from, take, calc, include, omit, infix, infixl, infixr, notation,
    macro, syntax, elab, set_option, initialize, deriving, opaque, abbrev,
    mutual, partial, scoped, termination_by, decreasing_by},
  % Sorts
  morekeywords=[2]{Sort, Type, Prop},
  % tactics, list taken from lean-syntax.txt (tactic keywords)
  morekeywords=[3]{intro, intros, apply, exact, refine, rfl, simp, simp_all,
    simpa, simp_rw, rw, rewrite, erewrite, ring, ring_nf, omega, decide,
    norm_num, norm_cast, induction, cases, rcases, obtain, constructor,
    ext, funext, conv, congr, field_simp, linarith, nlinarith, positivity,
    tauto, trivial, aesop, use, by_cases, by_contra, push_neg, contrapose,
    symm, trans, assumption, contradiction, exfalso, specialize,
    generalize, revert, clear, rename_i, subst, unfold, delta, change,
    convert, split_ifs, interval_cases, fin_cases, zify, gcongr,
    nth_rewrite, suffices, generalizing, at, only, native_decide},
  % warnings - constitution tier: must never survive into a final proof
  morekeywords=[4]{sorry, admit, sorryAx},
  literate=
    {α}{{\ensuremath{\alpha}}}1
    {β}{{\ensuremath{\beta}}}1
    {γ}{{\ensuremath{\gamma}}}1
    {δ}{{\ensuremath{\delta}}}1
    {ε}{{\ensuremath{\varepsilon}}}1
    {ζ}{{\ensuremath{\zeta}}}1
    {η}{{\ensuremath{\eta}}}1
    {θ}{{\ensuremath{\theta}}}1
    {ι}{{\ensuremath{\iota}}}1
    {κ}{{\ensuremath{\kappa}}}1
    {λ}{{\ensuremath{\lambda}}}1
    {μ}{{\ensuremath{\mu}}}1
    {ν}{{\ensuremath{\nu}}}1
    {ξ}{{\ensuremath{\xi}}}1
    {π}{{\ensuremath{\pi}}}1
    {ρ}{{\ensuremath{\rho}}}1
    {σ}{{\ensuremath{\sigma}}}1
    {τ}{{\ensuremath{\tau}}}1
    {υ}{{\ensuremath{\upsilon}}}1
    {φ}{{\ensuremath{\varphi}}}1
    {χ}{{\ensuremath{\chi}}}1
    {ψ}{{\ensuremath{\psi}}}1
    {ω}{{\ensuremath{\omega}}}1
    {Γ}{{\ensuremath{\Gamma}}}1
    {Δ}{{\ensuremath{\Delta}}}1
    {Θ}{{\ensuremath{\Theta}}}1
    {Λ}{{\ensuremath{\Lambda}}}1
    {Ξ}{{\ensuremath{\Xi}}}1
    {Π}{{\ensuremath{\Pi}}}1
    {Σ}{{\ensuremath{\Sigma}}}1
    {Φ}{{\ensuremath{\Phi}}}1
    {Ψ}{{\ensuremath{\Psi}}}1
    {Ω}{{\ensuremath{\Omega}}}1
    {ℵ}{{\ensuremath{\aleph}}}1
    {∀}{{\ensuremath{\forall}}}1
    {∃!}{{\ensuremath{\exists}!}}1
    {∃}{{\ensuremath{\exists}}}1
    {∈}{{\ensuremath{\in}}}1
    {∉}{{\ensuremath{\notin}}}1
    {∘}{{\ensuremath{\circ}}}1
    {∩}{{\ensuremath{\cap}}}1
    {∪}{{\ensuremath{\cup}}}1
    {⊆}{{\ensuremath{\subseteq}}}1
    {⊂}{{\ensuremath{\subset}}}1
    {⊇}{{\ensuremath{\supseteq}}}1
    {⊃}{{\ensuremath{\supset}}}1
    {⊄}{{\ensuremath{\nsubseteq}}}1
    {∅}{{\ensuremath{\emptyset}}}1
    {∧}{{\ensuremath{\wedge}}}1
    {∨}{{\ensuremath{\vee}}}1
    {¬}{{\ensuremath{\neg}}}1
    {⊤}{{\ensuremath{\top}}}1
    {⊥}{{\ensuremath{\bot}}}1
    {↔}{{\ensuremath{\leftrightarrow}}}1
    {→}{{\ensuremath{\rightarrow}}}1
    {←}{{\ensuremath{\leftarrow}}}1
    {↑}{{\ensuremath{\uparrow}}}1
    {⇒}{{\ensuremath{\Rightarrow}}}1
    {⇐}{{\ensuremath{\Leftarrow}}}1
    {⟹}{{\ensuremath{\Longrightarrow}}}1
    {↦}{{\ensuremath{\mapsto}}}1
    {≤}{{\ensuremath{\leq}}}1
    {≥}{{\ensuremath{\geq}}}1
    {≠}{{\ensuremath{\neq}}}1
    {≈}{{\ensuremath{\approx}}}1
    {≡}{{\ensuremath{\equiv}}}1
    {≃}{{\ensuremath{\simeq}}}1
    {≅}{{\ensuremath{\cong}}}1
    {×}{{\ensuremath{\times}}}1
    {·}{{\ensuremath{\cdot}}}1
    {±}{{\ensuremath{\pm}}}1
    {⊕}{{\ensuremath{\oplus}}}1
    {⊗}{{\ensuremath{\otimes}}}1
    {⋆}{{\ensuremath{\star}}}1
    {▸}{{\ensuremath{\blacktriangleright}}}1
    {⟨}{{\ensuremath{\langle}}}1
    {⟩}{{\ensuremath{\rangle}}}1
    {⦃}{{\ensuremath{\{\!\mid}}}1
    {⦄}{{\ensuremath{\mid\!\}}}}1
    {⟦}{{\ensuremath{[\![}}}1
    {⟧}{{\ensuremath{]\!]}}}1
    {⌊}{{\ensuremath{\lfloor}}}1
    {⌋}{{\ensuremath{\rfloor}}}1
    {⌈}{{\ensuremath{\lceil}}}1
    {⌉}{{\ensuremath{\rceil}}}1
    {√}{{\ensuremath{\surd}}}1
    {∣}{{\ensuremath{\mid}}}1
    {∤}{{\ensuremath{\nmid}}}1
    {∎}{{\ensuremath{\blacksquare}}}1
    {⋯}{{\ensuremath{\cdots}}}1
    {…}{{\ensuremath{\ldots}}}1
    {∑}{{\ensuremath{\sum}}}1
    {∏}{{\ensuremath{\prod}}}1
    {⁻¹}{{\ensuremath{^{-1}}}}1
    {¹}{{\ensuremath{^1}}}1
    {²}{{\ensuremath{^2}}}1
    {³}{{\ensuremath{^3}}}1
    {ⁿ}{{\ensuremath{^n}}}1
    {ᵐ}{{\ensuremath{^m}}}1
    {₀}{{\ensuremath{_0}}}1
    {₁}{{\ensuremath{_1}}}1
    {₂}{{\ensuremath{_2}}}1
    {₃}{{\ensuremath{_3}}}1
    {₄}{{\ensuremath{_4}}}1
    {₅}{{\ensuremath{_5}}}1
    {ₙ}{{\ensuremath{_n}}}1
    {ₘ}{{\ensuremath{_m}}}1
    {ₖ}{{\ensuremath{_k}}}1
    {ᵢ}{{\ensuremath{_i}}}1
    {ⱼ}{{\ensuremath{_j}}}1
    {ℤ}{{\ensuremath{\mathbb{Z}}}}1
    {ℕ}{{\ensuremath{\mathbb{N}}}}1
    {ℚ}{{\ensuremath{\mathbb{Q}}}}1
    {ℝ}{{\ensuremath{\mathbb{R}}}}1
    {ℂ}{{\ensuremath{\mathbb{C}}}}1
    {𝔽}{{\ensuremath{\mathbb{F}}}}1
    {𝒫}{{\ensuremath{\mathcal{P}}}}1
    {∗}{{\ensuremath{\ast}}}1
    {•}{{\ensuremath{\bullet}}}1,
  % Comments
  morecomment=[l]{--},
  morecomment=[s]{/-}{-/},
  morestring=[b]",
  stringstyle=\color{stringcolor},
  keepspaces=true,
  keywordstyle=[1]\color{keywordcolor}\bfseries,
  keywordstyle=[2]\color{sortcolor}\bfseries,
  keywordstyle=[3]\color{tacticcolor}\bfseries,
  keywordstyle=[4]\color{errorcolor}\bfseries\underline,
  escapeinside={(*@}{@*)},
  columns=fullflexible,
  basicstyle=\ttfamily\small,
  commentstyle=\color{commentcolor}\itshape,
  breaklines=true,
  showstringspaces=false,
  frame=single,
  framesep=2mm,
  xleftmargin=4mm,
  numbers=left,
  numberstyle=\tiny\color{commentcolor},
  numbersep=6pt,
}

% Inline Lean code in prose: \lean{Int.ModEq.pow}
\newcommand{\lean}[1]{\lstinline[language=lean,basicstyle=\ttfamily\normalsize]|#1|}

\lstset{language=lean}
% --- XeTeX rendering patch (backported from papers/cyl3-mod23/paper.tex,
%     extended with <= >= |- <-> emdash). Rendering-only; source bytes of
%     listings unchanged. Guarded so pdflatex keeps literate-table behaviour.
\ifdefined\XeTeXversion
  {\lccode`\~="2022 \lowercase{\global\catcode`~=\active \global\def~{\ensuremath{\bullet}}}}
  {\lccode`\~="2115 \lowercase{\global\catcode`~=\active \global\def~{\ensuremath{\mathbb{N}}}}}
  {\lccode`\~="2124 \lowercase{\global\catcode`~=\active \global\def~{\ensuremath{\mathbb{Z}}}}}
  {\lccode`\~="2190 \lowercase{\global\catcode`~=\active \global\def~{\ensuremath{\leftarrow}}}}
  {\lccode`\~="2192 \lowercase{\global\catcode`~=\active \global\def~{\ensuremath{\rightarrow}}}}
  {\lccode`\~="2194 \lowercase{\global\catcode`~=\active \global\def~{\ensuremath{\leftrightarrow}}}}
  {\lccode`\~="2200 \lowercase{\global\catcode`~=\active \global\def~{\ensuremath{\forall}}}}
  {\lccode`\~="2203 \lowercase{\global\catcode`~=\active \global\def~{\ensuremath{\exists}}}}
  {\lccode`\~="2227 \lowercase{\global\catcode`~=\active \global\def~{\ensuremath{\wedge}}}}
  {\lccode`\~="2228 \lowercase{\global\catcode`~=\active \global\def~{\ensuremath{\vee}}}}
  {\lccode`\~="2261 \lowercase{\global\catcode`~=\active \global\def~{\ensuremath{\equiv}}}}
  {\lccode`\~="2264 \lowercase{\global\catcode`~=\active \global\def~{\ensuremath{\leq}}}}
  {\lccode`\~="2265 \lowercase{\global\catcode`~=\active \global\def~{\ensuremath{\geq}}}}
  {\lccode`\~="22A2 \lowercase{\global\catcode`~=\active \global\def~{\ensuremath{\vdash}}}}
  {\lccode`\~="27E8 \lowercase{\global\catcode`~=\active \global\def~{\ensuremath{\langle}}}}
  {\lccode`\~="27E9 \lowercase{\global\catcode`~=\active \global\def~{\ensuremath{\rangle}}}}
\fi
\ifdefined\XeTeXversion
  \usepackage{xeCJK}
  \setCJKmainfont{SimSun}
\fi

\newtheorem{theorem}{Theorem}

\title{A mod-8 period-4 theorem and two companion identities for
recurrence-defined tatami counting sequences\\
(Lean 4 machine-checked claim of record)}
\author{Aurora0134\thanks{Contact: via the GitHub account Aurora0134; ORCID to
be added at Zenodo upload.}}
\date{September 27, 2026}

\begin{document}
\maketitle

\begin{abstract}
We record three theorems about two recurrence-defined integer sequences:
$a(n)$, given by the base values and the order-6 recurrence of OEIS entry
A180970, and $b(n)$, a corner-defect companion sequence satisfying the same
order-6 recurrence. For all $n \geq 4$, $a(n) \bmod 8$ depends only on
$n \bmod 4$, with pattern $(4,2,4,6)$; $b(n)$ is odd for all $n \geq 4$; and
$20\,b(n)$ equals a fixed linear combination of $a(n), \dots, a(n-5)$ for all
$n \geq 10$. All proofs are machine-checked by the Lean 4 kernel (0
\texttt{sorry}; axioms within \{propext, Quot.sound, Classical.choice\}). A
structured novelty search found no occupying record. The artifact package is
deposited on Zenodo. This note is a priority claim of record, not a full
paper.
\end{abstract}

\section{Introduction}

Let $a(n)$ be the integer sequence with base values
$a(0), \dots, a(8) = 1, 3, 13, 22, 44, 90, 196, 406, 852$ and the order-6
recurrence $a(n) = a(n-1) + 2a(n-2) + 2a(n-4) - a(n-5) - a(n-6)$ for
$n \geq 9$; these are exactly the base values and the recurrence signature
$(1,2,0,2,-1,-1)$ of OEIS entry A180970, the number of tatami tilings of a
$3 \times n$ grid with monomers allowed \cite{oeis-a180970}. The
identification of this recurrence-defined sequence with the tatami counting
sequence rests on the OEIS record and the Erickson--Ruskey literature
\cite{auspicious,ejc-monodimer,square-tatami} and is \emph{not} itself
formalized here. Let $b(n)$ (the corner-defect companion) be the integer
sequence with $b(0) := 0$ (a conventional value on which no theorem below
depends), machine-enumerated definition inputs
$b(1), \dots, b(9) = 2, 8, 24, 41, 85, 177, 381, 787, 1655$, and the same
order-6 recurrence for $n \geq 10$.

This note records three theorems about these two recurrence-defined
sequences, produced by AI4Math, a highly automated LLM-plus-kernel pipeline
in which the Lean 4 kernel is the sole arbiter of correctness. The present
text is a priority claim of record --- a short deposit note accompanying a
Zenodo artifact package; a fuller narrative paper is prepared separately
through the pipeline's paper lane. Every theorem stated below is a
\emph{complete proof} in the pipeline's grading: it compiles with 0
\texttt{sorry}, its axioms stay within the standard whitelist, and the
statement was frozen before proving.

\section{Statement}

The formal development defines the two sequences and the residue pattern
\lean{pat8} as follows (frozen snapshot; the definitions below are
byte-identical to the frozen statement file \texttt{01-statement.lean},
sha256 \texttt{59560f57...82b187f}).

\begin{lstlisting}[caption={Formal definitions of \texttt{a180970}, \texttt{pat8} and \texttt{bcorner} (verbatim excerpt, lines 12--54 of \texttt{final/01-proof-complete.lean}; byte-identical to the frozen statement file).},label={lst:defs}]
/-- A180970 as a recurrence-defined sequence: base values a(0..8) and the
order-6 recurrence a(n) = a(n-1) + 2 a(n-2) + 2 a(n-4) - a(n-5) - a(n-6)
for n ≥ 9 (OEIS signature (1,2,0,2,-1,-1)). -/
def a180970 : ℕ → ℤ
  | 0 => 1
  | 1 => 3
  | 2 => 13
  | 3 => 22
  | 4 => 44
  | 5 => 90
  | 6 => 196
  | 7 => 406
  | 8 => 852
  | n + 9 =>
      a180970 (n + 8) + 2 * a180970 (n + 7) + 2 * a180970 (n + 5)
        - a180970 (n + 4) - a180970 (n + 3)

/-- Residue pattern: a(n) % 8 = pat8 (n % 4). -/
def pat8 : ℕ → ℤ
  | 0 => 4
  | 1 => 2
  | 2 => 4
  | _ => 6

/-- Corner-defect sequence b_c(n), recurrence-defined version: tatami tilings of a
3×n grid with one corner cell removed.  Base values b(1..9) from exact enumeration
(machine probe, n ≤ 20); b(0) := 0 is a conventional value (a 3×0 grid has no corner
cell) — no theorem below depends on it.  For n ≥ 10 the same order-6 recurrence as
the base sequence holds: b(n) = b(n-1) + 2 b(n-2) + 2 b(n-4) - b(n-5) - b(n-6). -/
def bcorner : ℕ → ℤ
  | 0 => 0
  | 1 => 2
  | 2 => 8
  | 3 => 24
  | 4 => 41
  | 5 => 85
  | 6 => 177
  | 7 => 381
  | 8 => 787
  | 9 => 1655
  | n + 10 =>
      bcorner (n + 9) + 2 * bcorner (n + 8) + 2 * bcorner (n + 6)
        - bcorner (n + 5) - bcorner (n + 4)
\end{lstlisting}

% LEAN: tasks/20260926-beian-select :: tatami_mod8_period4  grade=完全证明
\begin{theorem}[mod-8 period 4]\label{thm:p1}
For every integer $n \geq 4$, the residue $a(n) \bmod 8$ depends only on
$n \bmod 4$: with $p(0) = 4$, $p(1) = 2$, $p(2) = 4$ and $p(3) = 6$, one has
$a(n) \bmod 8 = p(n \bmod 4)$.
\end{theorem}

% LEAN: tasks/20260926-beian-select :: bcorner_odd  grade=完全证明
\begin{theorem}[corner-defect parity]\label{thm:p2}
For every integer $n \geq 4$, $b(n)$ is odd.
\end{theorem}

% LEAN: tasks/20260926-beian-select :: bcorner_eq  grade=完全证明
\begin{theorem}[defect--base identity]\label{thm:p3}
For every integer $n \geq 10$,
\[
20\,b(n) \;=\; 4a(n) + 33a(n-1) - 9a(n-2) + 8a(n-3) + a(n-4) - 3a(n-5).
\]
\end{theorem}

\section{Proof}

The complete machine proofs are reproduced below verbatim from the frozen
final development. Listing~\ref{lst:p1} and Listing~\ref{lst:p2} are the full
proofs of Theorem~\ref{thm:p1} and Theorem~\ref{thm:p2};
Listing~\ref{lst:p3} shows the statement of Theorem~\ref{thm:p3} (its
97-line proof is included in full in the artifact package,
\texttt{proofs/01-proof-complete.lean}). All listings are excerpts of the
single file \texttt{01-proof-complete.lean} (207 lines, sha256
\texttt{b0f4a692...4789a1}). A narrative prose proof is not the duty of this
note; it belongs to the full paper lane.

\begin{lstlisting}[caption={Complete proof of Theorem~\ref{thm:p1} (verbatim excerpt, lines 56--84 of \texttt{final/01-proof-complete.lean}).},label={lst:p1}]
/-- **P1 (主攻)**: the mod-8 residue of A180970(n) depends only on n mod 4,
with pattern (4, 2, 4, 6) — equivalently a(n) % 8 = pat8 (n % 4) for all n ≥ 4.
Statement form verbatim from the r2 smoke-round kernel-accepted text. -/
theorem tatami_mod8_period4 (n r : ℕ) (hn : 4 ≤ n) (hr : n % 4 = r) :
    a180970 n % 8 = pat8 r := by
  induction n using Nat.strong_induction_on generalizing r with
  | _ n ih =>
    rcases lt_or_ge n 10 with h | h
    · interval_cases n
      all_goals norm_num at hr
      all_goals subst r
      all_goals norm_num [a180970, pat8]
    · obtain ⟨j, rfl⟩ : ∃ j, n = j + 9 := ⟨n - 9, by omega⟩
      have hj : 1 ≤ j := by omega
      have rec9 : ∀ m : ℕ, a180970 (m + 9)
          = a180970 (m + 8) + 2 * a180970 (m + 7) + 2 * a180970 (m + 5)
            - a180970 (m + 4) - a180970 (m + 3) := fun m => by
        simp only [a180970]
      rw [rec9 j]
      have h8 := ih (j + 8) (by omega) ((r + 3) % 4) (by omega) (by omega)
      have h7 := ih (j + 7) (by omega) ((r + 2) % 4) (by omega) (by omega)
      have h5 := ih (j + 5) (by omega) r (by omega) (by omega)
      have h4 := ih (j + 4) (by omega) ((r + 3) % 4) (by omega) (by omega)
      have h3 := ih (j + 3) (by omega) ((r + 2) % 4) (by omega) (by omega)
      have hml := Nat.mod_lt (j + 9) (show 0 < 4 by norm_num)
      have rl : r < 4 := by omega
      interval_cases r <;>
        simp only [Nat.reduceAdd, Nat.reduceMod, pat8] at h8 h7 h5 h4 h3 ⊢ <;>
        omega
\end{lstlisting}

\begin{lstlisting}[caption={Complete proof of Theorem~\ref{thm:p2} (verbatim excerpt, lines 86--104 of \texttt{final/01-proof-complete.lean}).},label={lst:p2}]
/-- **P2 (伴随)**: the corner-defect sequence b_c(n) is odd for all n ≥ 4. -/
theorem bcorner_odd (n : ℕ) (hn : 4 ≤ n) :
    bcorner n % 2 = 1 := by
  induction n using Nat.strong_induction_on with
  | _ n ih =>
    rcases lt_or_ge n 10 with h | h
    · interval_cases n <;> norm_num [bcorner]
    · obtain ⟨j, rfl⟩ : ∃ j, n = j + 10 := ⟨n - 10, by omega⟩
      have rec10 : ∀ m : ℕ, bcorner (m + 10)
          = bcorner (m + 9) + 2 * bcorner (m + 8) + 2 * bcorner (m + 6)
            - bcorner (m + 5) - bcorner (m + 4) := fun m => by
        simp only [bcorner]
      rw [rec10 j]
      have h9 := ih (j + 9) (by omega) (by omega)
      have h8 := ih (j + 8) (by omega) (by omega)
      have h6 := ih (j + 6) (by omega) (by omega)
      have h5 := ih (j + 5) (by omega) (by omega)
      have h4 := ih (j + 4) (by omega) (by omega)
      omega
\end{lstlisting}

\begin{lstlisting}[caption={Formal statement of Theorem~\ref{thm:p3} (verbatim excerpt, lines 106--110 of \texttt{final/01-proof-complete.lean}); proof (97 lines) in artifact package \texttt{proofs/}.},label={lst:p3}]
/-- **P3 (伴随)**: exact defect↔base linear identity for n ≥ 10:
20·b(n) = 4a(n) + 33a(n-1) - 9a(n-2) + 8a(n-3) + a(n-4) - 3a(n-5). -/
theorem bcorner_eq (n : ℕ) (hn : 10 ≤ n) :
    20 * bcorner n = 4 * a180970 n + 33 * a180970 (n - 1) - 9 * a180970 (n - 2)
      + 8 * a180970 (n - 3) + a180970 (n - 4) - 3 * a180970 (n - 5) := by
\end{lstlisting}

\section{Verification evidence}

\begin{itemize}
\item Toolchain: Lean v4.34.0 (\texttt{leanprover/lean4:v4.34.0})
\cite{lean4}, mathlib4 v4.34.0 (rev \texttt{5ed29652}) \cite{mathlib}.
\item Statement integrity: frozen statement
\texttt{formalized/01-statement.lean}, sha256
\texttt{59560f5735f2c8b58a31d631a4ed28806ca23292d195fa098231fb04682b187f};
final development \texttt{final/01-proof-complete.lean}, sha256
\texttt{b0f4a692990680e67817e0975484718132a6c54b9a1664598e44d9d0534789a1};
symmetric difference after stripping proof bodies: zero (59/59 lines
character-identical).
\item Kernel check: strict compile exit 0 with empty diagnostics; a
byte-level scan of the whole file (comments included) finds 0 occurrences of
\texttt{sorry}/\texttt{admit}. Axiom audit as recorded in the final-audit
record (\texttt{audit/final-01-audit.md}):
\begin{verbatim}
tatami_mod8_period4 : [propext, Classical.choice, Quot.sound]  (exactly the whitelist)
bcorner_odd         : [propext, Quot.sound]                    (subset)
bcorner_eq          : [propext, Quot.sound]                    (subset)
sorryAx             : absent
\end{verbatim}
\item Audit chain: gate-2 well-definedness compilation plus degeneracy
probing (24/24 bare-tactic probes fail, so no theorem is closable by a bare
decision procedure), statement freeze with symmetric-difference comparison,
gate-3 final strict compilation and axiom audit. Evidence originals are in
the artifact package \texttt{audit/}.
\end{itemize}

\section{Novelty evidence}

The pipeline's C9 novelty review adjudicated this result as \emph{no
occupying record found} (proposition layer). This is a statement about
search evidence, not a claim to new mathematics; the object A180970 itself
is an occupied OEIS record \cite{oeis-a180970}, and this note claims
properties of the recurrence-defined sequences, not the entry. Query
channels and query-string summary (full query log, including zero-hit
evidence files, in \texttt{audit/c9-record.md} of the package, copied
verbatim from the source records):

\begin{center}
\begin{tabular}{p{0.22\textwidth}p{0.47\textwidth}p{0.22\textwidth}}
\hline
channel & query summary & result \\
\hline
OEIS sequence layer & 9 strings: corner / short-edge / long-edge defect
sequences, each original and with 1 or 2 leading terms dropped & all null \\
OEIS sequence layer & positive control $22, 44, 90, 196, 406, 852$ &
hits A180970 (channel works) \\
OEIS record layer & full A180970 record re-fetched; case-insensitive scan
for \texttt{mod|period|congru} & no hit \\
literature layer & arXiv \texttt{all:"tatami"} (40 results);
\texttt{abs:"tatami" AND abs:"congruence"}; OpenAlex \texttt{tatami tilings}
& 0 relevant; no occupant \\
library layer & mathlib, compfiles, sequencelib, local rg for
\texttt{tatami} / \texttt{A180970} & zero hits \\
\hline
\end{tabular}
\end{center}

The exclusivity re-review (2026-09-27, independent re-query before deposit)
maintained the adjudication \emph{no occupying record found}; its recorded
residual gaps: CNKI not searched (login required), Semantic Scholar not
re-tested, the teaching-literature layer (textbooks and problem sheets) not
systematically swept, and the record covers the proposition layer only --- not
the combinatorial identification with tatami counting, which is layered as in
Section~\ref{sec:scope}.

\section{Scope and limitations}\label{sec:scope}

Grading boundary, verbatim from the final audit: the three theorems are
complete proofs, strictly as statements about the recurrence-defined integer
sequences; they do not cover any combinatorial counting meaning. Not
formalized here: (i) the identification $a(n) = \text{A180970}(n) = {}$the
tatami tiling count, which is backed by the OEIS record and the
Erickson--Ruskey literature \cite{oeis-a180970,auspicious,ejc-monodimer}
(and the square-grid closed forms of \cite{square-tatami} as background
only); and (ii) the reading of $b(n)$ as the number of tatami tilings of a
$3 \times n$ grid with one corner cell removed, which remains at the
computational-conjecture layer (exact machine enumeration probe for
$n \leq 20$ plus transfer-matrix general theory) and is not part of the
proved content. The base values $b(1), \dots, b(9)$ are machine-enumerated
definition inputs; $b(0) := 0$ is a conventional value that no theorem
depends on. This note contains no literature survey; it is a deposit record,
not a survey article.

\section*{Data availability}

Artifact package: Zenodo deposit (DOI to be reserved at upload; see
UPLOAD.md in the source repository) with a public mirror on GitHub (URL recorded in
claims/tatami-mod8-defect/card.md upon push). Note under CC BY 4.0, code
under MIT.

\section*{Use of generative AI}

(a) Stage mapping: candidate-proposition generation, formalization into Lean
statements, proof search, and text drafting (including this note) were
performed by AI (LLM) stages of the AI4Math pipeline; topic adjudication,
statement-semantics verification and gate sign-off were human. (b)
Model/node/budget (numbers copied verbatim from the source task ledger
\texttt{tasks/20260926-beian-select/budget.log}): a single resolved model
node, wire-id \texttt{anthropic/a6api-main/kimi-k3} (resolution level L2
relay), no cross-node switching; LLM sampling for the whole source task
totals llm-call $2/68$ (selection smoke $1$, tokens in/out $109/45$;
formalization roundtrip $1$, tokens $821/2418$); Lean compilations $22/36$
(plus $22/24$ on the selection smoke track). This claim-note stage used
llm-call $0$ of its budget $4$, and $8$ LaTeX compilations of its budget $6$
(counting the final rebuild that embedded this number and the author handle;
ledger \texttt{claims/tatami-mod8-defect/budget.log}). (c) Final
arbiter: All statements and proofs reported here were checked by the Lean 4
kernel: every proof compiles with no \texttt{sorry}, and \texttt{\#print
axioms} reports only \{propext, Quot.sound, Classical.choice\}. (d) The
human author reviewed and is responsible for all content; the AI system is
not an author.

\section*{Author contributions}

CRediT-style: the human author was responsible for topic selection and
adjudication, verification of statement semantics, gate sign-off, and the
decision to deposit; the AI4Math pipeline (an LLM system) performed
proposition generation, proof search, and text drafting. The AI system is
not listed as an author.

\begin{thebibliography}{9}

\bibitem{lean4}
L.~de Moura and S.~Ullrich.
\newblock The Lean 4 theorem prover and programming language.
\newblock \emph{CADE-28}, LNCS 12699, pp.~625--635, 2021.
\url{https://doi.org/10.1007/978-3-030-79876-5_37}

\bibitem{mathlib}
The mathlib Community.
\newblock The Lean mathematical library.
\newblock \emph{CPP 2020}, pp.~367--381, 2020.
\url{https://doi.org/10.1145/3372885.3373824}

\bibitem{auspicious}
A.~Erickson, F.~Ruskey, M.~Schurch, and J.~Woodcock.
\newblock Auspicious tatami mat arrangements.
\newblock \emph{arXiv:1103.3309 [math.CO]}, 2011.
\url{https://arxiv.org/abs/1103.3309}

\bibitem{ejc-monodimer}
A.~Erickson, F.~Ruskey, J.~Woodcock, and M.~Schurch.
\newblock Monomer-dimer tatami tilings of rectangular regions.
\newblock \emph{Electron. J. Combin.} 18(1), \#P109, 2011.
\url{https://doi.org/10.37236/596}

\bibitem{square-tatami}
A.~Erickson and M.~Schurch.
\newblock Monomer-dimer tatami tilings of square regions.
\newblock \emph{arXiv:1110.5103 [math.CO]}, 2011.
\url{https://arxiv.org/abs/1110.5103}

\bibitem{oeis-a180970}
The OEIS Foundation.
\newblock Entry A180970 (F.~Ruskey, 2010): number of tatami tilings of a
$3 \times n$ grid (with monomers allowed).
\newblock \emph{The On-Line Encyclopedia of Integer Sequences}.
\url{https://oeis.org/A180970}

\end{thebibliography}

\end{document}
